% GENERATED FILE. Edit canonical lessons, then run npm run generate:books.
% canonical-source-hash: fnv1a64:d3ac8c7c2c4bf056
% canonical-lessons: ZH-C271-jidong, ZH-C271-gandong, ZH-C271-chijing, ZH-C271-wuliao, ZH-C271-xinshi, ZH-R271-review-266-271

\chapter{Moods}
\label{ch:zh-moods}
\hlchaptermodality{271}

\begin{chapteropening}
\textbf{By the end of this chapter:} \emph{I can say I am worked up, moved, taken aback or bored, or that something is on my mind.}

Complete the last lesson of chapter 271: I can say I am worked up, moved, taken aback or bored, or that something is on my mind.
\end{chapteropening}

\section[jīdòng]{jīdòng — excited; worked up}

\label{lesson:ZH-C271-jidong}



\begin{quote}
Before the new one: say the Mandarin for to back someone up, then the Mandarin for to believe.
\end{quote}

\subsection*{You'll want to know: jīdòng}
\emph{jīdòng} — \textquotedblleft{}excited; worked up\textquotedblright{}.

\begin{quote}
\emph{tā hěn jīdòng} — He's very excited.
\end{quote}

Its characters come later; for now it is a word for the ear.

\begin{cousinweb}
\emph{jīdòng} is built from \emph{jī} and \emph{dòng} 'to move', both inherited from Middle Chinese, \emph{kek} and \emph{duwngX} (Guangyun, in Baxter's transcription), which glosses the first 'a rushing wave' and the second 'to stir'.
\end{cousinweb}

\begin{grammarlens}[title={What you've built}]
One more word for describing things.
\end{grammarlens}

\subsection*{Your turn}
\begin{itemize}
\raggedright
  \item \emph{Say it:} \emph{jīdòng}
  \item \emph{Say it:} \emph{jīdòng}, once more
  \item \emph{Say it:} \emph{tā hěn jīdòng} — He's very excited.
  \item \emph{Recall:} say the Mandarin for to back someone up, then the Mandarin for to believe, then say \emph{jīdòng} again
\end{itemize}

\begin{tcolorbox}[breakable,colback=teal!4,colframe=teal!35!black,title={Before you move on}]
What does jīdòng mean? (\textquotedblleft{}Excited; worked up\textquotedblright{}.) Say it once more.
\end{tcolorbox}
\section[gǎndòng]{gǎndòng — to move; to touch someone's heart}

\label{lesson:ZH-C271-gandong}



\begin{quote}
Before the new one: say the Mandarin for to believe, then the Mandarin for worked up.
\end{quote}

\subsection*{You'll want to know: gǎndòng}
\emph{gǎndòng} — \textquotedblleft{}to move; to touch someone's heart\textquotedblright{}.

\begin{quote}
\emph{zhège gùshi gǎndòngle dàjiā} — This story moved everyone.
\end{quote}

Its characters come later; for now it is a word for the ear.

\begin{cousinweb}
\emph{gǎndòng} is built from \emph{gǎn} and \emph{dòng}, both inherited from Middle Chinese, \emph{komX} and \emph{duwngX} (Guangyun, in Baxter's transcription), which glosses the first 'to move' and the second 'to stir'.
\end{cousinweb}

\begin{grammarlens}[title={What you've built}]
One more word for describing things.
\end{grammarlens}

\subsection*{Your turn}
\begin{itemize}
\raggedright
  \item \emph{Say it:} \emph{gǎndòng}
  \item \emph{Say it:} \emph{gǎndòng}, once more
  \item \emph{Say it:} \emph{zhège gùshi gǎndòngle dàjiā} — This story moved everyone.
  \item \emph{Recall:} say the Mandarin for to believe, then the Mandarin for worked up, then say \emph{gǎndòng} again
\end{itemize}

\begin{tcolorbox}[breakable,colback=teal!4,colframe=teal!35!black,title={Before you move on}]
What does gǎndòng mean? (\textquotedblleft{}To move; to touch someone's heart\textquotedblright{}.) Say it once more.
\end{tcolorbox}
\section[chījīng]{chījīng — surprised; taken aback}

\label{lesson:ZH-C271-chijing}



\begin{quote}
Before the new one: say the Mandarin for worked up, then the Mandarin for to touch someone's heart.
\end{quote}

\subsection*{You'll want to know: chījīng}
\emph{chījīng} — \textquotedblleft{}surprised; taken aback\textquotedblright{}.

\begin{quote}
\emph{wǒ hěn chījīng} — I'm surprised.
\end{quote}

Its characters come later; for now it is a word for the ear.

\begin{cousinweb}
\emph{chījīng} is built from \zh{吃} 'to eat' and \emph{jīng} 'a fright'. The second is inherited from Middle Chinese \emph{kjaeng} (Guangyun, in Baxter's transcription), glossed 'to fear' and, after an older dictionary, 'a horse taking fright'; how eating a fright came to mean surprise is not settled.
\end{cousinweb}

\begin{grammarlens}[title={What you've built}]
One more word for describing things.
\end{grammarlens}

\subsection*{Your turn}
\begin{itemize}
\raggedright
  \item \emph{Say it:} \emph{chījīng}
  \item \emph{Say it:} \emph{chījīng}, once more
  \item \emph{Say it:} \emph{wǒ hěn chījīng} — I'm surprised.
  \item \emph{Recall:} say the Mandarin for worked up, then the Mandarin for to touch someone's heart, then say \emph{chījīng} again
\end{itemize}

\begin{tcolorbox}[breakable,colback=teal!4,colframe=teal!35!black,title={Before you move on}]
What does chījīng mean? (\textquotedblleft{}Surprised; taken aback\textquotedblright{}.) Say it once more.
\end{tcolorbox}
\section[wúliáo]{wúliáo — bored; boring}

\label{lesson:ZH-C271-wuliao}



\begin{quote}
Before the new one: say the Mandarin for to touch someone's heart, then the Mandarin for taken aback.
\end{quote}

\subsection*{You'll want to know: wúliáo}
\emph{wúliáo} — \textquotedblleft{}bored; boring\textquotedblright{}.

\begin{quote}
\emph{zhège jiémù hěn wúliáo} — This show is boring.
\end{quote}

Its characters come later; for now it is a word for the ear.

\begin{cousinweb}
\emph{wúliáo} is built from \emph{wú} 'not to have' and \emph{liáo}, both inherited from Middle Chinese, \emph{mju} and \emph{lew} (Guangyun, in Baxter's transcription), which glosses the first 'to have or not have' and calls the second a particle. How the pair came to mean bored is not settled.
\end{cousinweb}

\begin{grammarlens}[title={What you've built}]
One more word for describing things.
\end{grammarlens}

\subsection*{Your turn}
\begin{itemize}
\raggedright
  \item \emph{Say it:} \emph{wúliáo}
  \item \emph{Say it:} \emph{wúliáo}, once more
  \item \emph{Say it:} \emph{zhège jiémù hěn wúliáo} — This show is boring.
  \item \emph{Recall:} say the Mandarin for to touch someone's heart, then the Mandarin for taken aback, then say \emph{wúliáo} again
\end{itemize}

\begin{tcolorbox}[breakable,colback=teal!4,colframe=teal!35!black,title={Before you move on}]
What does wúliáo mean? (\textquotedblleft{}Bored; boring\textquotedblright{}.) Say it once more.
\end{tcolorbox}
\section[\texorpdfstring{\zh{心事} (xīnshì)}{心事 (xīnshì)}]{\zh{心事} (xīnshì) — something on your mind; a worry}

\label{lesson:ZH-C271-xinshi}



\begin{quote}
Before the new one: say the Mandarin for taken aback, then the Mandarin for bored.
\end{quote}

\subsection*{You'll want to know: \zh{心事}}
\textbf{\zh{心事}} — \emph{xīnshì} — \textquotedblleft{}something on your mind; a worry\textquotedblright{}.

\begin{quote}
\textbf{\zh{他今天不笑有心事}} \emph{tā jīntiān bú xiào yǒu xīnshì} — He hasn't smiled today; something's on his mind.
\end{quote}

\begin{cousinweb}
\textbf{\zh{心事}} (\emph{xīnshì}) is built from \zh{心} 'the heart' and \zh{事} 'an affair', both inherited from Middle Chinese, \emph{sim} and \emph{dzriH} (Guangyun, in Baxter's transcription).
\end{cousinweb}

\begin{grammarlens}[title={What you've built}]
One more word for describing things.
\end{grammarlens}

\subsection*{Your turn}
\begin{itemize}
\raggedright
  \item \emph{Say it:} \emph{xīnshì}
  \item \emph{Say it:} \emph{xīnshì}, once more
  \item \emph{Say it:} \emph{tā jīntiān bú xiào yǒu xīnshì} — He hasn't smiled today; something's on his mind.
  \item \emph{Recall:} say the Mandarin for taken aback, then the Mandarin for bored, then say \emph{xīnshì} again
\end{itemize}

\begin{tcolorbox}[breakable,colback=teal!4,colframe=teal!35!black,title={Before you move on}]
What does \zh{心事} mean? (\textquotedblleft{}Something on your mind; a worry\textquotedblright{}.) Say it once more.
\end{tcolorbox}
\section[(dialogue)]{Looking back: from Grown-Ups and Children to Moods}

\label{lesson:ZH-R271-review-266-271}



\begin{quote}
Say the Mandarin for a grown-up and a baby.
\end{quote}

\subsection*{Your turn}
\begin{itemize}
\raggedright
  \item \emph{Say it:} a grown-up, a baby, a boy, a girl, middle age
  \item \emph{Say it:} an actor, a TV show, a celebrity, a joke, funny
  \item \emph{Say it:} to sign up, an oral exam, to pass an exam, the answer to a question, to explain
  \item \emph{Say it:} since childhood, before anything else, childhood, originally, a toy
  \item \emph{Say it:} to agree, to be against, to be in favour, to back someone up, to believe
  \item \emph{Say it:} worked up, to touch someone's heart, taken aback, bored, something on your mind
  \item \emph{Read it:} \zh{笑} \zh{先} — each one aloud, then a word that holds it
\end{itemize}

\begin{tcolorbox}[breakable,colback=teal!4,colframe=teal!35!black,title={Before you move on}]
A grown-up? (\textbf{\zh{大人}}, \emph{dàren}.) A baby? (\emph{bǎobao}.) A boy? (\textbf{\zh{男孩}}, \emph{nánhái}.) A girl? (\textbf{\zh{女孩}}, \emph{nǚhái}.) Middle age? (\textbf{\zh{中年}}, \emph{zhōngnián}.) An actor? (\emph{yǎnyuán}.) A TV show? (\emph{jiémù}.) A celebrity? (\textbf{\zh{明星}}, \emph{míngxīng}.) A joke? (\textbf{\zh{笑话}}, \emph{xiàohua}.) Funny? (\textbf{\zh{好笑}}, \emph{hǎoxiào}.) To sign up? (\emph{bàomíng}.) An oral exam? (\textbf{\zh{口试}}, \emph{kǒushì}.) To pass an exam? (\emph{jígé}.) The answer to a question? (\emph{dá'àn}.) To explain? (\emph{jiěshì}.) Since childhood? (\textbf{\zh{从小}}, \emph{cóngxiǎo}.) Before anything else? (\textbf{\zh{先}}, \emph{xiān}.) Childhood? (\emph{tóngnián}.) Originally? (\textbf{\zh{本来}}, \emph{běnlái}.) A toy? (\emph{wánjù}.) To agree? (\emph{tóngyì}.) To be against? (\emph{fǎnduì}.) To be in favour? (\emph{zànchéng}.) To back someone up? (\emph{zhīchí}.) To believe? (\emph{xiāngxìn}.) Worked up? (\emph{jīdòng}.) To touch someone's heart? (\emph{gǎndòng}.) Taken aback? (\emph{chījīng}.) Bored? (\emph{wúliáo}.) Something on your mind? (\textbf{\zh{心事}}, \emph{xīnshì}.)
\end{tcolorbox}
